\documentclass[../main]{subfiles}
\begin{document}

\chapter*{Exposición Integrador Máquinas Térmicas}
\img{images/16/caldera.jpeg}{Calderas}

Las máquinas térmicas son dispositivos que transforman energía térmica en trabajo mecánico, y su importancia radica en su uso extendido en diversas industrias, desde la generación de electricidad hasta la climatización. Las calderas, en particular, son esenciales en muchos procesos industriales, y en este capítulo se detallarán los aspectos más relevantes de su funcionamiento y aplicaciones.

    \begin{table}[h!]
        \centering
        \begin{tabular}{|p{5cm} | p{10cm} |} \hline
            Tema & Subtema \\ \hline \hline 
             Seguridad & Dispositivos y accesorios de seguridad en calderas: válvulas de seguridad, manómetros, interruptores de presión y sistemas de emergencia.\\ \hline 
             Calderas de Uso industrial & Estructura y características de calderas utilizadas en procesos industriales de gran escala, su capacidad de generación de vapor y eficiencia térmica. \\ \hline
             Condensadores y Quemadores & Sistemas de combustión eficientes, manejo de gases de escape, chimenea, caja de humos, y los diversos intercambiadores de calor.\\ \hline
             Máquinas Frigoríficas& El ciclo inverso de Carnot aplicado a la refrigeración industrial y doméstica.\\ \hline
             Centrales térmicas y equipamiento auxiliar para calderas& Precalentadores, desgasificadores, bombas de alta presión y otros dispositivos auxiliares para mejorar la eficiencia.\\ \hline 
        \end{tabular}
        \caption{Temas de Calderas y Máquinas Térmicas}
    \end{table}

\section{Contenido de la presentación.}
En esta sección se detallará el contenido que se abordará en la exposición sobre máquinas térmicas, cubriendo los principios de funcionamiento, ejemplos prácticos y aspectos técnicos esenciales para comprender el uso y mantenimiento de estas máquinas.

\begin{enumerate}
    \item \emph{Principio de funcionamiento:} Explicar en detalle cada una de las partes que componen a la Máquina Térmica, como la cámara de combustión, los intercambiadores de calor, y los sistemas de alimentación, incluyendo tanto las especificaciones técnicas de cada componente como un ejemplo concreto de aplicación industrial. Se debe prestar especial atención a los principios de termodinámica que gobiernan su operación.
    \item Ciclo de funcionamiento: Explicar con detalle cómo circula el calor a lo largo del ciclo, describiendo cada fase del proceso (compresión, expansión, y rechazo de calor), y en qué se convierte la energía transferida. Se incluirá una explicación sobre el principio físico subyacente que caracteriza su funcionamiento, como el ciclo de Carnot o Rankine, dependiendo del tipo de máquina.
    \item \emph{Fotos}: Buscar al menos 10 fotos de diferentes tipos de máquinas térmicas, calderas, sistemas de refrigeración y generadores de energía, tanto en aplicaciones industriales como domésticas.
    \item \emph{Historia:} Relatar la historia de la fabricación de estas máquinas, con especial énfasis en las primeras aplicaciones en la industria. Se incluirá información sobre el inventor o creador del primer modelo funcional, así como los primeros usos documentados de este tipo de máquina en la industria, tales como su introducción en la generación de energía eléctrica o procesos de calefacción a gran escala.
    \item \emph{Características Técnicas}: Detallar las características técnicas de las máquinas, como el consumo energético, la potencia generada, el tipo de energía utilizada (combustibles fósiles, biomasa, etc.), los parámetros de combustión, presión de funcionamiento, caudal de fluido, entre otros. También se abordarán los requisitos de eficiencia y normativa vigente en relación al uso de estas máquinas.
    \item \emph{Mantenimiento}: Explicar cómo se realiza el mantenimiento preventivo y correctivo de estas máquinas, detallando los procedimientos recomendados, la frecuencia de inspección de los componentes críticos, y las posibles fallas comunes que deben atenderse. Incluir un análisis sobre los costos asociados al mantenimiento y la vida útil de los equipos.
    \item \emph{Ventajas y desventajas:} Se analizarán las principales ventajas de utilizar máquinas térmicas en diversas aplicaciones, tales como la capacidad de generar grandes cantidades de energía de manera eficiente, su versatilidad en distintas industrias, y la posibilidad de utilizar diferentes tipos de combustibles. También se discutirán las desventajas, como los problemas relacionados con la emisión de gases contaminantes y los altos costos iniciales de instalación.
    \item Mercado actual:
    \begin{itemize}
        \item \emph{¿Dónde se compra?} Se investigarán los principales proveedores y fabricantes de estas máquinas, tanto a nivel nacional como internacional, y los canales de distribución disponibles para su adquisición.
        \item \emph{¿Cuál es el precio?} Proveer una lista de precios estimados para diferentes modelos comerciales de calderas y máquinas térmicas, incluyendo opciones de baja, media y alta capacidad. Se especificarán los factores que influyen en los costos, como el tamaño, la eficiencia, y el tipo de combustible.
        \item \emph{¿Qué características tiene cada modelo comercial?} Se compararán varias marcas y modelos comerciales, indicando las principales diferencias en cuanto a rendimiento, consumo energético, durabilidad, y tecnología utilizada.
    \end{itemize}
    \item \emph{Conclusiones:} Se presentará un resumen de las conclusiones generales después de la presentación, destacando los puntos más importantes tratados durante la clase. Se incluirá una reflexión sobre la importancia de las máquinas térmicas en la actualidad y su rol en el futuro de la industria energética.
    \item \emph{Preguntas:} Al final de la presentación, se abrirá un espacio para responder preguntas, donde los alumnos podrán aclarar cualquier duda respecto a los temas tratados.
    \item Subir archivo PDF con el informe detallado y la presentación en formato PowerPoint (PPT), a la plataforma classroom, asegurándose de que el archivo esté etiquetado con los nombres completos de todos los integrantes del equipo.
\end{enumerate}

\end{document}
